% Example "brusselator": the problem, the two species in space-time and the
% oscillation at one point, and the convergence and stability tests. Every
% number comes from data/, written by examples/brusselator/main.f90: run
% `fpm run --example brusselator` from the project root, then build this file
% with lualatex in examples/brusselator/.
\documentclass[11pt,a4paper]{article}
\usepackage[margin=2.2cm]{geometry}
\usepackage{amsmath,booktabs,array}
\usepackage{pgfplots}
\usepgfplotslibrary{groupplots,colormaps}
\pgfplotsset{compat=1.18}
\pagestyle{empty}

\colorlet{m-euler}{black}
\colorlet{m-order2}{orange!85!black}
\colorlet{m-order3}{green!50!black}
\colorlet{m-rk4}{blue!80!black}
\colorlet{m-dp54}{red!75!black}

\begin{document}
\section*{Example: the Brusselator, reaction and diffusion}

\subsection*{1. Problem}
From the SciML problem library (\texttt{prob\_ode\_brusselator\_1d}), with the
parameters of Hairer and Wanner (\emph{Solving ODEs II}, section IV.1): two
chemical species react and diffuse on a periodic interval,
\begin{equation*}
  u_t = A - (B + 1)\,u + u^2 v + \alpha u_{xx}, \qquad
  v_t = B u - u^2 v + \alpha v_{xx},
\end{equation*}
\begin{equation*}
  A = 1,\; B = 3,\; \alpha = \tfrac1{50}, \qquad u(x, 0) = 1 + \sin 2\pi x,
  \quad v(x, 0) = 3, \qquad 0 \le x < 1 \text{ periodic}.
\end{equation*}
Without diffusion the reaction has the fixed point $(A, B/A) = (1, 3)$,
unstable since $B > 1 + A^2$: the concentrations oscillate on a limit cycle.
Central differences on $n$ points ($\Delta x = 1/n$) turn diffusion into
eigenvalues $-\tfrac{4\alpha}{\Delta x^2}\sin^2(k\pi\Delta x)$, the largest
in size $4\alpha/\Delta x^2$: refining the grid makes the system stiff, while
the reaction stays mild. No closed form; the reference is \texttt{dp54} at a
tiny step.
\begin{center}
\small
% Written by examples/brusselator/main.f90: do not edit.
\begin{tabular}{l>{\raggedright\arraybackslash}p{12.4cm}}
\toprule
parameters & $A = 1$, $B = 3$, $\alpha = 1/50$ \\
initial condition & $u = 1 + \sin 2\pi x$, $v = 3$; periodic on $[0, 1)$ \\
space & central differences, $n$ points, $\Delta x = 1/n$; $2n$ unknowns \\
convergence & $n = 64$, error at $t = 10$ against \texttt{dp54} with $\Delta t = 10^{-5}$; $\Delta t = 0.004, \dots, 0.004/64$; order from the two finest levels above $10^{-12}$, kept if within 0.3 of the next \\
stability & $n = 16, 32, 64, 128, 256, 512$; largest $\Delta t$ keeping every value finite and below 100 up to $t = 2$ and for at least 2000 steps (initial value plus $10^{-6}(-1)^j$ in $u$); bisection in $[0.3, 2]$ times $h^*\Delta x^2/(4\alpha)$ \\
figures & \texttt{rk4}, $n = 64$, $\Delta t = 10^{-3}$, $0 \le t \le 20$ \\
methods & \texttt{rk3ws}, \texttt{rk4}, \texttt{dp54} \\
interface & plain arrays, \texttt{rk\_array\_scheme\_t} \\
\bottomrule
\end{tabular}

\end{center}

\newpage
\subsection*{2. Solutions}
\begin{center}
\begin{tikzpicture}
  \begin{groupplot}[
      group style={group size=2 by 1, horizontal sep=0.6cm, y descriptions at=edge left}, % two maps
      width=6.2cm, height=11.2cm,                          % tall maps
      view={0}{90}, xmin=0, xmax=1, ymin=0, ymax=20,       % x across, t up
      xlabel={$x$}, ylabel={$t$},                          % space-time
      colormap/viridis, enlargelimits=false, axis on top,  % the maps fill the axes
      colorbar horizontal,                                 % one colour bar under each map
    ]
    \nextgroupplot[title={(a) $u(x, t)$},
      colorbar style={at={(0.5,-0.1)}, anchor=north, xlabel={$u$}, height=0.2cm}]
    \addplot3[surf, shader=interp, mesh/ordering=x varies] table[x=x, y=t, z=u] {data/spacetime.dat}; % species u
    \addplot[white, thin, densely dashed, domain=0:20] ({0.5}, {x});                       % where (c) looks
    \nextgroupplot[title={(b) $v(x, t)$},
      colorbar style={at={(0.5,-0.1)}, anchor=north, xlabel={$v$}, height=0.2cm}]
    \addplot3[surf, shader=interp, mesh/ordering=x varies] table[x=x, y=t, z=v] {data/spacetime.dat}; % species v
    \addplot[white, thin, densely dashed, domain=0:20] ({0.5}, {x});                       % where (c) looks
  \end{groupplot}
  % ---- (c) the oscillation at x = 1/2 ----
  \begin{axis}[
      at={($(group c2r1.north east)+(1.4cm,0cm)$)}, anchor=north west, % right of the maps
      width=5.2cm, height=5.4cm,                           % panel
      xmin=0, xmax=20, ymin=0, ymax=5,                     % concentrations in time
      xlabel={$t$}, title={(c) at $x = \tfrac12$},          % panel title
      axis background/.style={fill=black!2},               % very light background
      legend style={font=\tiny, at={(0.02,0.98)}, anchor=north west, fill=white, fill opacity=0.85, text opacity=1, draw=black!30},
    ]
    \addplot[blue!70!black, thick] table[x=t, y=u] {data/probe.dat};                 % u(1/2, t)
    \addlegendentry{$u$}
    \addplot[orange!85!black, thick] table[x=t, y=v] {data/probe.dat};               % v(1/2, t)
    \addlegendentry{$v$}
  \end{axis}
  % ---- (d) the same in the (u, v) plane ----
  \begin{axis}[
      at={($(group c2r1.north east)+(1.4cm,-7.0cm)$)}, anchor=north west, % below (c)
      width=5.2cm, height=5.4cm,                           % panel
      xlabel={$u$}, ylabel={$v$}, title={(d) $(u, v)$ at $x = \tfrac12$},
      axis background/.style={fill=black!2},               % very light background
      colormap/viridis, point meta min=0, point meta max=20, % colour = time
    ]
    \addplot[mesh, thick, point meta=\thisrow{t}] table[x=u, y=v] {data/probe.dat};   % the orbit, coloured by time
    \addplot[only marks, mark=+, mark size=4pt, very thick, black] coordinates {(1,3)}; % fixed point (A, B/A)
  \end{axis}
\end{tikzpicture}
\end{center}
(a, b) The two species in space and time, \texttt{rk4} on 64 points. The
initial bump in $u$ is smoothed out by diffusion within the first time unit;
then the whole interval oscillates almost in unison, with a period of about
$7.1$: $u$ stays low while $v$ builds up, until $u$ spikes in a short burst
that uses up $v$.
(c) The oscillation at $x = \tfrac12$ (white dashed line in the maps). (d) The
same in the $(u, v)$ plane, coloured by time: the solution spirals away from
the unstable fixed point $(1, 3)$ (cross) onto the limit cycle of the
reaction.

\newpage
\subsection*{3. Tests: convergence and the stability limit}

\paragraph{Prediction.} (i) The problem is nonlinear: each method converges
at its general order (\texttt{rk3ws}: 2). (ii) On fine grids the fastest
eigenvalue is the diffusive one, $\approx -4\alpha/\Delta x^2$, so the stable
step should approach $h^*\Delta x^2/(4\alpha)$; on coarse grids the reaction,
whose Jacobian adds eigenvalues of order one, shifts it.

\paragraph{Result.}
\begin{center}
\small
% Written by examples/brusselator/main.f90: do not edit.
\begin{tabular}{lccccc}
\toprule
& & & & \multicolumn{2}{c}{stable $\Delta t$ / $h^*\Delta x^2/(4\alpha)$} \\
\cmidrule(lr){5-6}
name & $h^*$ & order & observed & $n = 16$ & $n = 512$ \\
\midrule
\texttt{rk3ws} & 2.5127 & 2 & 2.00 & 0.9485 & 1.0002 \\
\texttt{rk4} & 2.7853 & 4 & 4.02 & 0.9446 & 1.0002 \\
\texttt{dp54} & 3.3066 & 5 & --- & 0.9417 & 1.0001 \\
\bottomrule
\end{tabular}

\end{center}

\begin{center}
\begin{tikzpicture}
  \begin{groupplot}[
      group style={group size=2 by 1, horizontal sep=1.8cm}, % two panels
      width=7.4cm, height=6.6cm,                           % size of each panel
      grid=major, grid style={black!8},                    % faint grid
      axis background/.style={fill=black!2},               % very light background
      every axis plot/.append style={thick, mark size=1.3pt}, % same weight for all curves
    ]
    % ---- (i) convergence ----
    \nextgroupplot[
      xmode=log, ymode=log, title={(i) error at $t = 10$, $n = 64$}, xlabel={$\Delta t$},
      legend to name=brlegend, legend columns=7,           % legend below the panels
      legend style={font=\small, /tikz/every even column/.append style={column sep=0.25cm}},
    ]
    \addplot[m-order3, dashed, mark=triangle] table[x=dt, y=error] {data/error_rk3ws.dat};
    \addlegendentry{\texttt{rk3ws}}
    \addplot[m-rk4, mark=diamond*]            table[x=dt, y=error] {data/error_rk4.dat};
    \addlegendentry{\texttt{rk4}}
    \addplot[m-dp54, mark=pentagon*]          table[x=dt, y=error] {data/error_dp54.dat};
    \addlegendentry{\texttt{dp54}}
    % ---- (ii) stable step relative to the diffusion prediction ----
    \nextgroupplot[
      xmode=log, title={(ii) stable $\Delta t$ / $h^*\Delta x^2/(4\alpha)$}, xlabel={$\Delta x$},
      ymin=0.9, ymax=1.03, xtick={0.002,0.004,0.008,0.016,0.031,0.0625},
      xticklabels={$\frac1{512}$,$\frac1{256}$,$\frac1{128}$,$\frac1{64}$,$\frac1{32}$,$\frac1{16}$},
    ]
    \addplot[black!40, thin, dashed, domain=0.0015:0.08] {1};                                   % the prediction
    \addplot[m-order3, dashed, mark=triangle] table[x=dx, y expr=\thisrow{measured}/\thisrow{predicted}] {data/stable_rk3ws.dat};
    \addplot[m-rk4, mark=diamond*]            table[x=dx, y expr=\thisrow{measured}/\thisrow{predicted}] {data/stable_rk4.dat};
    \addplot[m-dp54, mark=pentagon*]          table[x=dx, y expr=\thisrow{measured}/\thisrow{predicted}] {data/stable_dp54.dat};
  \end{groupplot}
  \node[below=0.2cm] at (current bounding box.south)        % legend centred under the panels
    {\pgfplotslegendfromname{brlegend}};
\end{tikzpicture}
\end{center}
(i) Every method converges at its general order (\texttt{rk3ws} at 2 on this
nonlinear problem); \texttt{dp54} is at round-off already at the largest step
(---). (ii) The measured stable step divided by the diffusion prediction:
on the coarsest grid the reaction lowers it by $4$--$7\,\%$; from 64 points on
it agrees with $h^*\Delta x^2/(4\alpha)$ within $0.4\,\%$ for every method,
and within $0.03\,\%$ on 512 points. At $n = 512$ that is a step of about
$10^{-4}$ for \texttt{rk4}, on a problem whose oscillation takes seven time
units: diffusion, not the chemistry, sets the explicit step.
\end{document}
